\section{TIM-MARS Formulation}

TIM-MARS is implemented as a selected-target memory layer above a tracker. At frame $t$, the tracker provides a candidate set
\begin{equation}
\mathcal{C}_t = \{c_t^1,\ldots,c_t^N\},
\end{equation}
where each candidate is represented as
\begin{equation}
c_t^i = (b_t^i, q_t^i, id_t^i, e_t^i).
\end{equation}
Here, $b_t^i$ is the bounding box, $q_t^i$ is the tracker confidence, $id_t^i$ is the tracker identity, and $e_t^i$ is the MARS appearance descriptor when available.

The memory state is
\begin{equation}
m_t = (b_t^\star, id_t^\star, e_t^\star, s_t, h_t),
\end{equation}
where $b_t^\star$ is the remembered selected-target box, $id_t^\star$ is the remembered tracker identity, $e_t^\star$ is the appearance memory, $s_t$ is the finite memory state, and $h_t$ stores hysteresis counters.

\subsection{Base Candidate Score}

Candidates below a minimum tracker confidence are discarded:
\begin{equation}
\mathcal{C}_t' =
\{c_t^i \in \mathcal{C}_t \mid q_t^i \geq \tau_q\}.
\end{equation}

For each remaining candidate, TIM-MARS computes five terms:
\begin{align}
I_i &= \mathrm{IoU}(b_t^i,b_{t-1}^\star),\\
D_i &= \exp\left[-\frac{1}{2}\left(\frac{d(b_t^i,b_{t-1}^\star)}{\sigma_d}\right)^2\right],\\
Z_i &= \exp\left[-\frac{|\log(A(b_t^i)/A(b_{t-1}^\star))|}{\sigma_z}\right],\\
Q_i &= \mathrm{clip}(q_t^i,0,1),\\
B_i &=
\begin{cases}
1, & id_t^i = id_{t-1}^\star,\\
0, & \mathrm{otherwise}.
\end{cases}
\end{align}

The base score is
\begin{equation}
S_i^{0}
=
w_I I_i
+
w_D D_i
+
w_Z Z_i
+
w_Q Q_i
+
w_B B_i.
\end{equation}

The base implementation weights are
\begin{equation}
(w_I,w_D,w_Z,w_Q,w_B)
=
(0.34,0.26,0.18,0.14,0.08).
\end{equation}

\subsection{Ambiguity Test}

Let $c_t^{(1)}$ and $c_t^{(2)}$ be the highest and second-highest candidates under $S_i^0$. The base ambiguity margin is
\begin{equation}
\Delta_t^0 = S_{(1)}^0 - S_{(2)}^0.
\end{equation}

The candidate set is treated as ambiguous when
\begin{equation}
\Delta_t^0 < \tau_\Delta,
\quad
\tau_\Delta = 0.07.
\end{equation}

\subsection{Appearance Gating}

MARS appearance is not applied unconditionally. First, appearance must be available for both the candidate and the memory. Second, geometry must allow appearance use:
\begin{equation}
\Gamma_i =
[I_i > 0]
\lor
[D_i \geq 0.25 \land Z_i \geq 0.35].
\end{equation}

The raw appearance similarity is cosine similarity:
\begin{equation}
A_i =
\frac{e_t^i \cdot e_{t-1}^\star}
{\lVert e_t^i \rVert_2 \lVert e_{t-1}^\star \rVert_2}.
\end{equation}

If appearance is enabled, available, geometry-allowed, and passes the similarity gate,
\begin{equation}
A_i \geq \tau_A,
\end{equation}
then the final score is
\begin{equation}
S_i = \mathrm{clip}(S_i^0 + w_A A_i,0,1).
\end{equation}
Otherwise,
\begin{equation}
S_i = S_i^0.
\end{equation}

In the hard re-entry evaluation,
\begin{equation}
w_A = 0.30,
\quad
\tau_A = 0.35.
\end{equation}

Figure~\ref{fig:scoring_block} summarises the candidate scoring path.

\begin{figure}[t]
\centering
\begin{tikzpicture}[
    node distance=4mm,
    block/.style={draw, rounded corners, align=center, minimum width=25mm, minimum height=7mm, font=\scriptsize},
    small/.style={draw, rounded corners, align=center, minimum width=19mm, minimum height=6mm, font=\scriptsize},
    arrow/.style={-Latex, thick}
]
\node[block] (cand) {candidate\\$c_t^i$};
\node[small, below left=of cand] (geom) {geometry\\$I_i,D_i,Z_i$};
\node[small, below=of cand] (conf) {tracker\\$Q_i,B_i$};
\node[small, below right=of cand] (app) {MARS\\$A_i$};

\node[block, below=12mm of conf] (score) {score\\$S_i=S_i^0+w_AA_i$};
\node[block, below=of score] (accept) {threshold\\and margin gates};
\node[block, below=of accept] (state) {LOCKED / REACQUIRED\\or UNCERTAIN / LOST};

\draw[arrow] (cand) -- (geom);
\draw[arrow] (cand) -- (conf);
\draw[arrow] (cand) -- (app);
\draw[arrow] (geom) -- (score);
\draw[arrow] (conf) -- (score);
\draw[arrow] (app) -- (score);
\draw[arrow] (score) -- (accept);
\draw[arrow] (accept) -- (state);
\end{tikzpicture}
\caption{TIM-MARS candidate scoring and rejection flow. Geometry and tracker continuity define the base score, while MARS appearance is added only when available and gated.}
\label{fig:scoring_block}
\end{figure}


\subsection{Acceptance Threshold}

Let $c_t^\star$ be the highest-scoring candidate under $S_i$. The acceptance threshold depends on the previous memory state:
\begin{equation}
\tau_s =
\begin{cases}
\tau_L, & s_{t-1} = \mathrm{LOST},\\
\tau_K, & \mathrm{otherwise},
\end{cases}
\end{equation}
with
\begin{equation}
\tau_K = 0.52,
\quad
\tau_L = 0.60.
\end{equation}

If the candidate keeps the same tracker identity as the memory, the threshold is relaxed:
\begin{equation}
\tau_s' =
\max(0,\tau_s - \delta_{id}),
\quad
\delta_{id}=0.08.
\end{equation}

A candidate is rejected before acceptance if
\begin{equation}
S_t^\star < \tau_s'.
\end{equation}

\subsection{Conservative Appearance Rejection}

After a candidate passes the score threshold, TIM-MARS can still reject it using conservative appearance filtering.

Let $A^\star$ be the selected candidate raw appearance similarity, and let $A^{(2)}$ be the second-highest raw appearance similarity among candidates for which appearance was used and geometry allowed appearance. The appearance margin is
\begin{equation}
\Delta_A = A^\star - A^{(2)}.
\end{equation}

With conservative filtering enabled, the selected candidate is rejected if
\begin{equation}
A^\star < \tau_C
\lor
\Delta_A < \tau_M.
\end{equation}

For the hard re-entry evaluation,
\begin{equation}
\tau_C = 0.65,
\quad
\tau_M = 0.25.
\end{equation}

If appearance was not used, the missing-appearance rejection is
\begin{equation}
R_\emptyset =
\neg \mathrm{appearance\_used}(c_t^\star)
\land
\rho_A,
\end{equation}
where
\begin{equation}
\rho_A =
\begin{cases}
1, & \texttt{appearance\_conservative\_require\_appearance=true},\\
0, & \texttt{appearance\_conservative\_require\_appearance=false}.
\end{cases}
\end{equation}

The evaluated policy uses
\begin{equation}
\rho_A = 0.
\end{equation}
Therefore, missing appearance alone does not force a LOST state.

\subsection{State and Output}

The memory state is updated as
\begin{equation}
s_t =
\begin{cases}
\mathrm{REACQUIRED}, & s_{t-1}\in\{\mathrm{LOST},\mathrm{UNCERTAIN}\} \land \mathrm{accepted},\\
\mathrm{LOCKED}, & \mathrm{accepted},\\
\mathrm{UNCERTAIN}, & \mathrm{temporarily\ rejected},\\
\mathrm{LOST}, & \mathrm{rejected\ for\ too\ long}.
\end{cases}
\end{equation}

The control-facing output is
\begin{equation}
y_t =
\begin{cases}
b_t^\star, & s_t\in\{\mathrm{LOCKED},\mathrm{REACQUIRED}\},\\
\emptyset, & s_t\in\{\mathrm{UNCERTAIN},\mathrm{LOST},\mathrm{NO\_TARGET}\}.
\end{cases}
\end{equation}

This implements the control-safety rule: ambiguous identity evidence should produce no target rather than an unsafe wrong target.
